\section{From heights to finiteness}\label{sec:counting}

We now turn the height estimate into a bound on degrees.  The rational
objects to be counted are ordered pairs of integral Hodge vectors in the
rank-five variation.  Their two orthogonality equations cut the period
quadric in a conic.  A semialgebraic family of such pairs with one real
parameter therefore sweeps out an algebraic subset of dimension at most
two, whereas full monodromy forces the period image to be Zariski dense
in the three-dimensional quadric.

Throughout this section the curve, family, and all auxiliary covers are
fixed over a number field $K$.  We choose the ample height on $\bar S$
from a fixed projective embedding defined over $K$, and write
\[
 h(t)=\max\{2,h_{\bar S}(t)\},\qquad
 d(t)=\max\{2,[K(t):K]\}.
\]
This choice is invariant under $\Gal(\Qbar/K)$; changing to it preserves
the preceding estimates by the height comparisons in
\Cref{lem:finite-covers}.  Constants in this section depend on the fixed
data.

\subsection{Flat coordinates at all complex embeddings}

The Rosati bound in \Cref{prop:small-vectors} controls the Hodge norms
of our vectors.  Counting requires their numerical sizes in integral flat
bases.  Near a puncture, passing between these two bounds also requires
control of the dual Hodge norm at each complex embedding.

Write $b=\pair{\cdot}{\cdot}$ for the trace pairing of
\eqref{eq:quadratic-system}.  In a flat trivialization of
$\mathbb V_{\Z}$, let $V_{\Z}$ be the corresponding fixed lattice and
$V_{\C}=V_{\Z}\otimes\C$.  The period quadric of
\eqref{eq:period-quadric} is
\[
 Q_V=\{[w]\in\mathbb P(V_{\C}):b(w,w)=0\},
 \qquad \dim_{\C}Q_V=3.
\]
The period line is the line corresponding to $H^{2,0}$, with the Tate
twist understood.  A real vector of type $(0,0)$ is orthogonal to this
line.  In particular the two vectors of \Cref{prop:small-vectors} satisfy
these equations, and their Gram matrix for $b$ is positive definite by
\Cref{lem:qm-plane}.

Here the counting height is the \emph{multiplicative rational coordinate
height}.  For $y=(y_1,\ldots,y_m)\in\Q^m$, with
$y_i=a_i/b_i$ in lowest terms and $b_i>0$, put
\[
 H_{\mathrm{rat}}(y)=\max_i\max\{\abs{a_i},b_i\}.
\]
For an integral tuple it is $\max\{1,\norm{y}_{\infty}\}$.  This height
is distinct from the absolute logarithmic algebraic heights used in the
arithmetic argument.

\begin{lemma}[Flat-coordinate bound]\label{lem:flat-coordinate-bound}
There is a finite collection of simply connected complex charts covering
$S(\C)$, each equipped with an integral flat basis of $\mathbb V_{\Z}$, whose
period maps are definable in $\R_{\mathrm{an},\exp}$.  There are constants
$c,M>0$ such that, for every QM point $t\in S(\Qbar)$, the ordered pair
of integral vectors supplied by \Cref{prop:small-vectors} has flat
coordinate tuple $y\in\Z^{10}$ satisfying
\begin{equation}\label{eq:flat-coordinate-height}
 H_{\mathrm{rat}}(y)\le c\bigl(d(t)h(t)\bigr)^M
\end{equation}
in each of these charts containing $t$.  The same bound holds for their
conjugate vectors at every conjugate of $t$ over $K$.
\end{lemma}

\begin{proof}
We first describe the charts and their metric estimates, including the
punctures of $S$.  A compact subset away from the finitely many points of
$\bar S\setminus S$ has a finite cover by relatively compact simply
connected coordinate disks.  Shrink these disks inside slightly larger
disks on which the flat trivializations and the period maps are
holomorphic.  Their real and imaginary parts are restricted analytic
functions, hence definable.

At a puncture choose a rational function $u$ on $\bar S$ which is a local
parameter there.  Enlarge $K$ once to define these finitely many choices.
The local monodromy is quasi-unipotent.  On a fixed local cover
$u=w^e$ it is unipotent, say $\exp N$, with $N$ nilpotent.  The
one-variable nilpotent orbit theorem
\cite[Theorem~4.9]{Schmid1973} gives, in the compact dual, the expression
\begin{equation}\label{eq:sector-period}
 \Phi(u)=
 \exp\left(\frac{\log w}{2\pi i}N\right)\Psi(w),
\end{equation}
where $\Psi$ extends holomorphically across $w=0$.  Cover a sufficiently
small punctured $u$-disk by finitely many open angular sectors of width
less than $2\pi$, and choose a branch of $w$ and of $\log w$ on each.
The exponential in \eqref{eq:sector-period} is a polynomial in $\log w$.
The branch of $w$ is algebraic, its argument is bounded, and
$\log\abs{w}$ is definable in $\R_{\mathrm{an},\exp}$.  After shrinking
the radius, $\Psi$ is restricted analytic in an affine chart about
$\Psi(0)$.  Thus the projective map \eqref{eq:sector-period} is
definable.  These local covers are used only to describe functions on
the sectors; no point-dependent global cover is introduced.

Write $\norm{\cdot}_{H,u}$ for the Hodge norm.  Schmid's abstract
one-variable norm estimate
\cite[Theorem~6.6$'$]{Schmid1973}, applied both to $\mathbb V$ and to its
dual, bounds the norms of each fixed flat basis and dual basis by a
power of $1+\log(1/\abs{u})$, uniformly on each sector.  Indeed every
fixed flat vector belongs to some step of the monodromy weight
filtration, and that theorem supplies such a power; only finitely many
basis vectors occur.  Passing from $w$ to $u=w^e$ merely changes the
constants.  Consequently, with $e_1,\ldots,e_5$ an integral flat basis
and $e_1^*,\ldots,e_5^*$ its dual, we have fixed $c_0,r>0$ such that
\begin{equation}\label{eq:sector-metric-dual}
 \max_i\bigl\{\norm{e_i}_{H,u},
                  \norm{e_i^*}_{H,u}^{\vee}\bigr\}
 \le c_0\bigl(1+\log(1/\abs{u})\bigr)^r.
\end{equation}
On the compact charts the corresponding bound is constant.  In
particular \eqref{eq:sector-metric-dual} controls the metric and its
inverse in flat coordinates.

For completeness, the logarithm in this bound can be controlled at an
individual complex embedding, not merely on average.  Let
$F=K(t)$ and let $\sigma:F\hookrightarrow\C$ fix the chosen embedding
of $K$.  Whenever $\sigma(t)$ belongs to the sector, $u(t)\ne0$, and
the absolute height identity $\ha{u(t)^{-1}}=\ha{u(t)}$ gives
\begin{align}
 \log^+\frac1{\abs{\sigma(u(t))}}
 &\le [F:\Q]\ha{u(t)}\notag\\
 &\le c_1d(t)h(t).\label{eq:individual-embedding-log}
\end{align}
The first inequality follows by retaining that archimedean summand in
the height of $u(t)^{-1}$, using the usual nonsquared absolute value.
For the second, the fixed rational function $u$ defines a morphism
$\bar S\to\mathbb P^1$, whose height is bounded above by a constant
times $1+h_{\bar S}$.  This also explains why algebraic local
parameters were chosen.

Let $v$ be either of the two chosen endomorphism vectors, at any such
embedding.  If $v$ realizes the symmetric endomorphism $\alpha$, the Weil
operator fixes its type $(0,0)$, and its rationality gives
\[
 \norm{v}_H^2=b(v,v)=\tfrac14\Tr(\alpha^2)
              =\tfrac14 q_{\lambda_t}(\alpha).
\]
The bound in
\Cref{prop:small-vectors} therefore gives
$\norm{v}_{H}\le c_2(d(t)h(t))^{r_2}$.  If
$v=\sum_i n_i e_i$, the dual norm inequality is
\[
 \abs{n_i}=\abs{e_i^*(v)}
 \le\norm{e_i^*}_{H}^{\vee}\norm{v}_{H}.
\]
Combining this with \eqref{eq:sector-metric-dual} and
\eqref{eq:individual-embedding-log} proves a polynomial bound on the
\emph{numerical sizes} of all ten integers.  This is precisely
\eqref{eq:flat-coordinate-height}.

The vectors are realizations of actual integral endomorphisms.
Conjugation preserves integrality, independence, and the rational
Rosati trace norms.  Although labeling the endomorphisms may require an
extension of $K(t)$ of bounded relative degree, each embedding of
$K(t)$ extends to that field.  Conjugating the labeled pair along any
such extension supplies the asserted pair at the conjugate point.
Neither the degree estimate for $u(t)$ nor the number of conjugate
points is diminished by this labeling extension.
\end{proof}

\subsection{Counting distinct curve arguments}

We use the following precise specialization of the counting theorem.
The structure $\R_{\mathrm{an},\exp}$ containing our period charts is
o-minimal by \cite[Corollary~5.13]{DMM1994}.
Let $Z\subset\R^m\times\R^n$ be definable in an o-minimal expansion of
the real field, and write $\pi_1,\pi_2$ for its two projections.  For
every $\eta>0$ there is $c_Z(\eta)$ such that, if $T\ge1$ and
\[
 \Sigma\subset
 \{(y,z)\in Z:y\in\Q^m,\ H_{\mathrm{rat}}(y)\le T\},
 \qquad \#\pi_2(\Sigma)>c_Z(\eta)T^\eta,
\]
then there is a continuous definable path
$\beta:[0,1]\to Z$ for which $\pi_1\beta$ is semialgebraic and real
analytic on $(0,1)$, while $\pi_2\beta$ is nonconstant.  This is
\cite[Corollary~7.2, with $k=1$ and no family parameter]{HabeggerPila2016}.
The corollary counts distinct second projections and imposes no
finite-fiber hypothesis on $Z\to\R^m$.  A constant first projection is
allowed, and causes no difficulty below.

\begin{proposition}\label{prop:counting-finiteness}
For fixed $A,\kappa>0$, the set of QM points $t\in S(\Qbar)$ satisfying
\begin{equation}\label{eq:counting-height-hypothesis}
 h(t)\le A d(t)^\kappa
\end{equation}
is finite.
\end{proposition}

\begin{proof}
The set in question is stable under $\Gal(\Qbar/K)$.  If its degrees
were unbounded, take a point $t$ of sufficiently large degree
$d=[K(t):K]$.  Its full orbit has $d$ distinct points on $S(\C)$, all
satisfying \eqref{eq:counting-height-hypothesis}.  Assign each conjugate
to one chart of \Cref{lem:flat-coordinate-bound}.  If there are $J$
charts, one chart contains at least $d/J$ of these distinct points.
For each of them choose a conjugate of the labeled pair at $t$.
Their integral coordinate tuples have
\begin{equation}\label{eq:counting-polynomial-T}
 H_{\mathrm{rat}}(y)\le T=c_3d^B
\end{equation}
for fixed $c_3,B>0$.

In the selected chart write $z\in D\subset\R^2$ for the curve
coordinate and $\Phi(z)\in Q_V$ for the period line.  Identify
the flat real space with $\R^5$.  Consider the fixed definable set
\begin{equation}\label{eq:counting-incidence}
 Z=\left\{((v_1,v_2),z)\in(\R^5)^2\times D:
 \begin{array}{l}
   (b(v_i,v_j))_{i,j=1}^2\text{ is positive definite},\\
   b(v_1,\Phi(z))=b(v_2,\Phi(z))=0
 \end{array}\right\}.
\end{equation}
Orthogonality to a projective line means orthogonality to any nonzero
representative; it is a well-defined algebraic incidence condition.
The coordinate $z$ is injective on the chart.  The tuples just chosen
thus form a set $\Sigma\subset Z$ with at least $d/J$ distinct second
projections and rational first coordinates bounded as in
\eqref{eq:counting-polynomial-T}.

Choose $\eta>0$ with $B\eta<1$.  The finitely many chart incidences
have a common counting constant, obtained by taking their maximum.
For large enough $d$ we have
\[
 d/J>c_Z(\eta)(c_3d^B)^\eta.
\]
The quoted corollary therefore gives a path
$\beta(s)=(\gamma(s),z(s))$ in $Z$, with $\gamma$ semialgebraic and
$z$ nonconstant.  It remains to show that this path would force the
period image into a proper algebraic subset of $Q_V$.  The first
projection supplies at most one algebraic parameter, and for each
nondegenerate pair the orthogonal period lines form a conic.

Let $\mathcal A\subset(\C^5)^2$ be the complex Zariski closure of
$\gamma((0,1))$.  A semialgebraic set of real dimension at most one
has real Zariski closure of dimension at most one, and complexification
preserves that dimension.  Hence
$\dim_{\C}\mathcal A\le1$.  Equivalently, after a finite
semialgebraic decomposition, the arc has Nash coordinate functions,
which are algebraic over its real parameter; their algebraic closure
has transcendence degree at most one.

Crucially, impose nondegeneracy before closing the orthogonality
incidence.  Let
\[
 \mathcal U=
 \{(v_1,v_2):\det(b(v_i,v_j))_{i,j=1}^2\ne0\}
 \subset(\C^5)^2,
\]
and define the locally closed algebraic set
\[
 \mathcal I=
 \{(v_1,v_2,[w])\in(\mathcal A\cap\mathcal U)\times Q_V:
       b(v_1,w)=b(v_2,w)=0\}.
\]
For each $(v_1,v_2)\in\mathcal A\cap\mathcal U$, their span is a
nondegenerate two-space.  Its perpendicular is a nondegenerate
three-space, whose isotropic lines form a smooth projective conic.
Every fiber of $\mathcal I\to\mathcal A\cap\mathcal U$ therefore
has dimension one, so
\[
 \dim_{\C}\mathcal I\le2.
\]
Let $X\subset Q_V$ be the Zariski closure of its image under
the second projection.  The dimension of an image and its closure
cannot exceed the dimension of its source, whence
\begin{equation}\label{eq:counting-dimension-two}
 \dim_{\C}X\le2.
\end{equation}
This construction closes only the incidence over
$\mathcal A\cap\mathcal U$; it does not add arbitrary fibers over
degenerate pairs.  Since the pairs on $\gamma$ have positive definite
Gram matrix, they lie in $\mathcal U$, and consequently
$\Phi(z(s))\in X$ for all $s\in(0,1)$.

The continuous nonconstant path $z(s)$ contains infinitely many distinct
points accumulating inside the chart.  Indeed one may restrict it to a
closed interior interval on which it remains nonconstant.  The inverse
image $\Phi^{-1}(X)$ is a complex analytic subset of a connected
one-dimensional chart.  The identity theorem therefore makes it the
entire chart.  Continuing the defining homogeneous polynomial
identities along $S$ gives
\[
 \widetilde\Phi(\widetilde S)\subset X,
\]
where $\widetilde S$ is the connected universal cover and the flat
trivialization defining $\widetilde\Phi$ extends the one in the chosen
chart.

Fix a period line $L$ in this image.  Equivariance places its entire
monodromy orbit $\Gamma L$ in $X$.  By \Cref{prop:monodromy},
$\Gamma$ is Zariski dense in $\SO(V_{\C},b)$, and this group acts
transitively on the nonsingular isotropic quadric $Q_V$.
The orbit $\Gamma L$ is therefore Zariski dense in $Q_V$.
This contradicts \eqref{eq:counting-dimension-two} and
$\dim Q_V=3$.

Degrees in \eqref{eq:counting-height-hypothesis} are thus bounded.
That inequality also bounds their ample heights, and Northcott's
theorem on the fixed projective curve $\bar S$ proves finiteness.
\end{proof}

\subsection{Descent to the coarse curve}

\begin{proof}[Proof of \Cref{thm:main}]
Suppose the stated locus on $C$ were infinite.  By
\Cref{lem:finite-covers}, normalize $C$ and take a connected component
of a fixed fine-level cover containing infinitely many lifted target
points.  Removing the finitely many points needed for the smooth
family and the algebraic charts leaves infinitely many.  Choose an
interior QM center and then the fixed auxiliary cover of
\Cref{prop:marking}.  All these choices, including their fields of
definition, are made once.

For every remaining target to which the construction is applied,
\Cref{thm:height} supplies a selected lift on this fixed cover with
$h(t)\le A d(t)^\kappa$, for fixed $A,\kappa$.  The degrees of a point
and its lift are comparable by the fixed degree of the finite map, and
the ample heights are comparable by the pullback of ample divisors,
as in \Cref{lem:finite-covers}.  Thus the constants are uniform as the
quaternion algebra and the target vary.

Apply \Cref{prop:counting-finiteness} on the fixed covering curve.
Its counting argument uses the full orbit over the fixed field $K$:
the height inequality for a selected lift is invariant under that
orbit, and conjugation preserves its two integral Hodge vectors and
their Rosati bounds.  A conjugate marking need not satisfy the
incidence imposed relative to the chosen center.  That incidence was
used to prove the height inequality; the incidence
\eqref{eq:counting-incidence} depends only on the target's Hodge
vectors.  Hence all these selected lifts belong to the finite set of
\Cref{prop:counting-finiteness}, whose final step uses Northcott's
theorem.  Their images on $C$ form a finite set.  Adding back the
finitely many removed points contradicts the assumed infinitude.
This proves the asserted finiteness for full geometric endomorphism
algebra an indefinite quaternion division algebra over $\Q$, with no
restriction on its discriminant or order.
\end{proof}
